% =========================================================================================
% Chapter 9 -- Conclusion and Perspectives
% =========================================================================================
\section{Conclusion and Perspectives}
\label{sec:conclusion}

\subsection{Summary of Contributions}
\label{subsec:conclusion-contributions}

The individual model architectures used here are standard and are not claimed as
contributions. EfficientNet-B0, U-Net and gradient-boosted trees are established tools applied
in established ways, and a reader should assume that any competent practitioner would have
reached comparable per-module figures on the same data.

What was built is the layer above them.

\paragraph{A shared output contract across heterogeneous models.} A segmentation network, two
classifiers and a tree ensemble write through one validated schema, so a reasoning layer can
consume all four without knowing which produced what. The contract is enforced rather than
conventional: four validation rules reject reports that would be ambiguous downstream, each rule
written against a failure that actually occurred.

\paragraph{Calibrated confidence carried end to end.} Each module's confidence is calibrated by
a method suited to its own miscalibration and travels with the finding to the point of use,
instead of being computed for a metrics table and discarded at the module boundary. Where a
calibrator is absent the output says so, and the reasoning layer discounts it.

\paragraph{An explicit representation of absent evidence and model scope.} The system
distinguishes a finding detected, a finding screened for and not seen, and a finding never
assessed --- three states that a simpler design collapses into two. Each module declares what it
cannot exclude, and those declarations reach the escalation policy rather than remaining in
documentation. The practical consequence is that an all-negative scan escalates instead of
reassuring.

\paragraph{A deterministic escalation policy independent of the language model.} The
safety-critical decision is computed from the structured state before the model is consulted and
is unchanged by anything the model does, including failing entirely. This is demonstrated rather
than asserted: 339 tests exercise it with no model present, and the decision was verified
identical across three backends including one that raises on every call.

\paragraph{An evidence interface that makes fabrication unrepresentable.} The reasoning layer
does not ask the model to avoid inventing evidence; it removes the ability. Citations are
identifiers drawn from an enumerated list of facts the state actually holds, so there is no
identifier for an invented observation, none for a test never performed, and none for a sentence
lifted from the reference corpus. Section~\ref{subsec:reasoning-ab} reports what that changed.

\paragraph{Decision support that does not pass through the model.} Severity, critical alerts,
recommended examinations and scenario routing are computed from the structured state by the
module described in Section~\ref{subsec:reasoning-decision-support}, which imports nothing
capable of reaching a language model. Severity is monotone, so a reassuring triage tier never
lowers it beside a critical alert. Every numeric cutoff sits in one versioned configuration
file and every alert names the bound it fired on, which makes an alert checkable rather than
merely plausible.

\paragraph{A refusal, implemented.} Therapeutic suggestions are required by the specification
and are gated rather than generated:

\begin{keystatement}
Therapeutic recommendations were deliberately gated behind retrieval of an authoritative,
citable protocol. In the absence of an appropriate protocol source, the system produces no
therapeutic recommendation rather than relying on the language model's parametric knowledge.
\end{keystatement}

\paragraph{An automated clinical report.} One standardised document assembling the encounter,
the reasoning and the decision support, with an archived copy carrying a content hash so that
what a clinician was shown can be recovered later and proven unaltered. Within it, a modality
never assessed occupies its own section and a withheld differential is rendered as withheld
--- both because absence inferred from silence is indistinguishable from a normal result.

\paragraph{A methodological contribution.} Case-level evaluation was established for three
datasets that did not provide it. Each required reconstructing the statistical unit before any
model could be honestly evaluated, and in two of the three the reconstruction changed the
reported figures materially --- downward.

\medskip
Taken together, the position this work supports is a narrow one and is worth stating in the
form it can be defended:

\begin{keystatement}
The complete MVP pipeline is implemented and technically verified. The system provides
multimodal POCUS-based decision support with deterministic safety controls,
retrieval-augmented reasoning, structured recommendations and automated reporting. Clinical
validation against expert ground truth remains outside the current internship scope.
\end{keystatement}

Every word of that is checkable. \emph{Implemented} and \emph{technically verified} are claims
about software, evidenced by 339 tests. \emph{Deterministic safety controls} names a property
demonstrated by running the pipeline against a failing model backend. And the final sentence
states what the work does not establish, in the same breath as what it does, because the two
are not separable here.

\subsection{Achieved Objectives}
\label{subsec:conclusion-objectives}

The objectives of Section~\ref{subsec:objectives} are assessed below against the evidence of
Chapters~\ref{sec:perception} to~\ref{sec:results}. Several are partial, and they are marked
partial: the limitations chapter immediately precedes this one and a reader who has just read it
will not be persuaded by a clean scorecard.

\subsubsection{Business Objectives}

\paragraph{Reduce time-to-diagnosis --- not demonstrated.} The system produces an escalation
decision in under two seconds and a differential within two minutes, which is fast enough to be
used during an examination. Whether it shortens time-to-diagnosis is a claim about clinical
workflow that requires clinical evaluation, and none was performed. What can be said is that
nothing in the measured runtime prevents it.

\paragraph{Reduce diagnostic error for less experienced operators --- not demonstrated.} This
objective requires comparing clinicians using the system against clinicians without it. No such
comparison exists in this work, and the per-module accuracy figures do not substitute for one.

\paragraph{Make uncertainty visible rather than implicit --- achieved.} This is the objective
the work most directly meets, and the one that survives the limitations. Every finding carries a
calibrated confidence; absent tests are recorded as absent and cannot be cited as evidence;
every module declares what it cannot exclude; and every abnormal value reaches the clinician
whether or not the reasoning model referred to it. Each of these is enforced in code and
verified by tests rather than left to convention.

\paragraph{Preserve the clinician's authority over the decision --- achieved.} The system
recommends next steps and does not prescribe treatment; a validator rejects output that
instructs treatment. It escalates rather than resolving disagreements it lacks the information
to settle, and it withholds a differential it cannot ground rather than presenting an unreliable
one. The clinician sees the evidence, the escalation and the reasoning, and decides.

\subsubsection{Data Science Objectives}

\paragraph{Computer vision models returning calibrated confidence --- achieved.} Three modules,
each calibrated, each reporting a probability rather than a bare score. Their accuracy is
moderate and is reported as such.

\paragraph{An explainability layer relating attention to anatomical regions --- partial.}
Grad-CAM was implemented for all three imaging modules and its output was measured rather than
inspected by eye. What was established is that attention is concentrated and content-dependent.
What was \emph{not} established is that it falls on the anatomically correct structure, which is
what the objective asked for and which requires expert annotation this project does not have.

\paragraph{A structured clinical state recording what was not measured --- achieved.} This is
the component the rest of the architecture rests on, and it does what the objective specifies.

\paragraph{A reasoning layer producing a grounded differential --- partial.} A ranked
differential is produced and every citation is grounded in the state by construction. But the
model reasons from part of the available evidence rather than all of it, phrases recommendations
more strongly than decision support should, and has never been evaluated against clinical ground
truth. Grounding was achieved; reasoning quality was not established.

\paragraph{An escalation policy from confidence, severity and completeness --- achieved.} Seven
deterministic triggers, computed before the model runs, tested by 11 dedicated tests within the
180.

\subsection{Current Limitations}
\label{subsec:conclusion-limitations}

Section~\ref{subsec:results-limitations} states these in full and they are not repeated here.
In summary: the case-level datasets are small, the images are curated public collections rather
than consecutive clinical series, only the triage model was evaluated across independent
sources, and the reasoning benchmark comprises five synthetic cases.

One sentence carries more weight than the rest. \textbf{This system has not been validated
against clinical outcomes and is not a diagnostic device.} Nothing in this work was measured
against what happened to a patient or against what a clinician concluded, and no figure reported
anywhere in it should be read as evidence of clinical utility.

\subsection{Difficulties Encountered}
\label{subsec:conclusion-difficulties}

Five obstacles shaped the method rather than merely delaying it.

\paragraph{Intermittent GPU availability.} Training ran on free Colab sessions with no
guarantee of allocation and a runtime that could be reclaimed mid-experiment. This made each
training run expensive enough that speculative experimentation was not affordable, which forced
a habit of estimating a candidate change before spending a run on it. The clearest instance is
the gallbladder taxonomy selection of Section~\ref{subsec:gallbladder-module}, where collapsing
an existing confusion matrix arithmetically predicted the accuracy each candidate grouping would
reach, and one configuration was trained instead of several. The constraint produced a better
method than abundance would have.

\paragraph{Datasets whose statistical unit was undocumented.} Three of the four datasets provide
frame-level or image-level labels with no explicit case identifier, and none states that
splitting on them would leak. Reconstructing case identity --- from directory structure, from
metadata, and for the lung dataset from a filename heuristic --- was a precondition for any
honest figure, and it lowered the reported results in two modules. The earlier, higher numbers
were not wrong through carelessness; they were wrong because the unit of analysis had been
inherited from the dataset's layout rather than chosen.

\paragraph{Taxonomies that did not match the images.} A gallbladder class whose contents did not
correspond to its label was found only by sampling one frame per case across the whole folder.
The first inspection, from a handful of frames, had judged it consistent. The lesson generalised
beyond that dataset and is stated as a rule in Section~\ref{subsec:gallbladder-module}: frames
from one case are not a sample of a class.

\paragraph{Hypotheses contradicted by measurement.} The cardiac generalisation gap was first
attributed to a labelling-convention mismatch, on evidence from a small number of images.
Connected-component analysis over the full external dataset disproved it. The hypothesis was
abandoned and the reversal is reported in Section~\ref{subsec:cardiac-module} rather than
quietly dropped, because a method that only records the hypotheses that survived gives a reader
no way to judge it.

\paragraph{Results that lived only inside a notebook session.} The lung module's decision
thresholds and probability calibrators were computed during training and, for the calibrators,
never written to a file. Everything downstream consequently fell back to a default threshold of
$0.5$, which for this model's compressed output range meant reporting every study as negative
--- a failure with no error message and no visible symptom, described in
Section~\ref{subsubsec:agent-extraction}. The thresholds were recoverable because the training
run had also written a results file; the calibrators were not, and remain lost.

The lesson is narrow and worth stating exactly: an artefact that exists only in the memory of a
notebook session does not exist. It also generalises, because the defect was invisible until the
prediction code became importable and a test asserted a property of it. Research code that
cannot be imported cannot be tested, and code that is not tested is not known to work --- a
statement that sounds obvious and was nonetheless the single most expensive omission in this
project.

\subsection{Future Work}
\label{subsec:conclusion-future}

\subsubsection{Additional organs}

Vascular imaging for deep vein thrombosis was deliberately excluded from the scope of this
internship, and the FAST examination was defined but not implemented. Both already return
\texttt{status: "not\_supported"} through the ordinary contract, which means adding them
requires a module and no change to anything that consumes one. The reasoning layer, the
escalation policy and the schema would be untouched.

The reference corpus already contains sourced passages for the FAST views, so the knowledge
those modules would need to be reasoned about is in place before the modules themselves.

\subsubsection{Video analysis}

Every imaging module here operates on frames. The lung module samples eight frames from a clip
and aggregates them by median, which is aggregation over still images rather than analysis of
motion.

This is a hard bound rather than a refinement. Lung sliding is defined by movement of the
pleural line between frames; the lung point is defined by the transition between sliding and
absent sliding at a location; cardiac wall motion abnormality is a temporal pattern. None of
these can be represented by a model that sees frames independently, however much data it is
given. Pneumothorax detection in particular --- absent from this system for want of positive
cases --- depends on exactly the temporal signs a frame-based model cannot see.

\subsubsection{Improved calibration and external validation}

The evidence these results most lack is external. Only the triage model was evaluated across
independent sources, and the 17-point spread it showed between them is the best available
indication of how much the imaging figures might move on data from another institution.

What is needed is larger case-level datasets, prospective collection rather than curated
archives, and validation on a cohort from a different institution than the one that produced
the training data. Calibration would also improve: the current calibrators are fitted and scored
on the same out-of-fold predictions, so the reported expected calibration errors are optimistic,
and a properly held-out calibration set would give an honest figure.

\subsubsection{Simulation and the digital twin}

The escalation route currently ends at a recommendation to examine further. The intended
continuation, and the reason the escalation output is named \texttt{simulation}, is an
environment in which a candidate decision could be explored against a patient model before being
committed to.

\textbf{No part of this was built during the internship.} It is stated here as a perspective
because the architecture was arranged to accommodate it --- escalation is a routing decision
with a named destination, and that destination is currently unimplemented --- and not because
any of it exists.

\subsubsection{Clinical evaluation}

The step that would make this system meaningful is the one it does not have. Every figure in
this report measures whether a model scores well or whether a safety mechanism behaves as
specified. None measures whether a clinician using the system reaches a better decision, or a
faster one, than a clinician without it.

That evaluation would be a study rather than an experiment: clinicians of varying experience,
real presentations, and outcomes measured against what actually happened to the patient. Its
endpoints would be time-to-diagnosis and diagnostic error --- the two business objectives of
Section~\ref{subsec:objectives} that this work explicitly could not demonstrate. Until it is
performed, the honest description of what has been built is a prototype whose safety properties
are verified and whose clinical value is unknown.
